\section{Referencias Bibliográficas}

\begin{thebibliography}{99}

\bibitem{herrera2025lambert}
Herrera Herrera, A. A., Méndez López, D. A., \& Cisneros Díaz, I. A. (2025). La función Lambert W: teoría, propiedades y aplicaciones. \textit{Ciencias Matemáticas}, 39(1), 11--19. DOI: \url{https://doi.org/10.5281/zenodo.17443639}. Recuperado a partir de \url{https://revistas.uh.cu/rcm/article/view/10777}.

\bibitem{pensando2021lambert}
Pensando Numéricamente. (2021). \textit{FUNCIÓN W de LAMBERT: Origen, propiedades, dominio y aplicaciones | Pensando Numéricamente} [Video]. YouTube. \url{https://www.youtube.com/watch?v=HU5_qbUqKBU}

\bibitem{POLJAK1995249}
S.~Poljak y F.~Rendl.
\newblock Solving the max-cut problem using eigenvalues.
\newblock \emph{Discrete Applied Mathematics}, 62(1):249--278, 1995.

\bibitem{bondy2008graph}
J.~A. Bondy y U.~S.~R. Murty.
\newblock \emph{Graph Theory}.
\newblock Graduate Texts in Mathematics, vol.~244.
\newblock Springer, London, 2008.

\bibitem{jungnickel2005graphs}
D.~Jungnickel.
\newblock \emph{Graphs, Networks and Algorithms}.
\newblock Editado por M.~Bronstein, A.~M. Cohen, H.~Cohen, D.~Eisenbud y B.~Sturmfels.
\newblock Algorithms and Computation in Mathematics, vol.~5.
\newblock Springer, Berlin, Heidelberg, 2005.

\bibitem{GODDARD2013839}
W.~Goddard y M.~A. Henning.
\newblock Independent domination in graphs: A survey and recent results.
\newblock \emph{Discrete Mathematics}, 313(7):839--854, 2013.

\bibitem{kleinberg2005}
Kleinberg, J., \& Tardos, É. (2005). \textit{Algorithm Design} (1st ed.). Pearson/Addison-Wesley. (Específicamente Capítulo 7: Network Flow).

\bibitem{clrs2022}
Cormen, T. H., Leiserson, C. E., Rivest, R. L., \& Stein, C. (2022). \textit{Introduction to Algorithms} (4th ed.). Cambridge, MA: MIT Press.

\end{thebibliography}